\chapter{L'impianto di valutazione}\label{chapter:valutazione}

Ogni numero di questa tesi viene da un file di un archivio, prodotto da codice
committato, e ogni misura che conta ha le sue previsioni scritte prima. Questo
capitolo descrive le regole, le collezioni, le query e gli strumenti; le regole
vengono prima, perché sono loro a dire quanto ci si può fidare del resto.

\section{Principi: ipotesi prima, interruttori, baseline congelato, archivio con provenienza}\label{sec:principi}

\paragraph{Le ipotesi prima della misura.} Ogni modifica che cambia l'ordine dei
risultati ha una voce nel diario di Koskidex, e ogni esperimento un README, con
una sezione \emph{Prima di misurare} committata prima di lanciare la
valutazione, spesso prima di scrivere il codice. Le previsioni sono numeriche,
con una soglia, e dicono cosa vorrebbe dire sbagliarle. A misura fatta si
aggiunge l'esito, previsione per previsione, senza toccare quello che era scritto
prima. Quello che si decide dopo aver visto un numero, un'analisi in più o una
strategia nuova, sta in una sezione che lo dichiara. È la risposta alla domanda
che in una valutazione arriva sempre: le metriche e le soglie sono state scelte
dopo aver visto i risultati? L'appendice~\ref{chapter:appendice-ipotesi} le raccoglie
tutte, e molte sono state smentite.

\paragraph{Gli interruttori.} Ogni correzione sta dietro un'impostazione, con il
comportamento di partenza come valore predefinito. Le impostazioni si salvano con
l'indice, e un indice salvato prima di una correzione non cambia comportamento da
solo. Così ogni misura confronta due configurazioni dello stesso codice, e i
default si decidono alla fine, tutti insieme.

\paragraph{Il baseline congelato.} Il comportamento di Koskidex del 23 settembre
2026 è fissato in un test, \texttt{TestBaselineRankingIsFrozen}, che confronta gli
ordinamenti di un corpus di prova con quelli registrati. Deve passare a ogni
commit, con le impostazioni predefinite e senza che il test venga toccato: è la
prova che nessuna correzione ha cambiato il punto di partenza di nascosto.

\paragraph{L'archivio.} Le esecuzioni vanno in un archivio, nel repository della
tesi, con regole scritte prima di riempirlo:
\begin{itemize}
  \item un file per esecuzione, con l'ora UTC nel nome, mai sovrascritto né
        modificato a mano;
  \item ogni file registra da dove viene: il commit del codice, se l'albero aveva
        modifiche non committate, il corpus con la sua impronta, ogni
        impostazione. Per la tesi valgono solo i file da alberi puliti;
  \item i tempi stanno separati dalle metriche, perché cambiano a ogni esecuzione
        e le metriche no;
  \item solo identificativi, mai titoli: gli atti contengono nomi di persone, e
        il corpus resta fuori dal repository;
  \item solo esecuzioni vere: le prove fatte mentre si scrive codice si lanciano
        con l'archivio spento, e gli strumenti, se l'archivio non è impostato,
        lo dicono invece di tacere.
\end{itemize}
Le eccezioni stanno in un registro: dodici file del primo giorno prodotti da
codice non ancora committato e rifatti, un file rinominato, tre ridotti da 12 MB a
1 MB conservando i primi mille risultati, e così via. Nessuno è stato tolto in
silenzio.

\paragraph{Le guardie.} Una previsione può essere una guardia: una quantità che la
modifica non deve muovere. La prima scritta, «BM25 non cambia il Recall@100», è
scattata perché era posta male (sezione~\ref{sec:difetto1}); da lì il valutatore
registra per ogni query il numero dei candidati prima del taglio, che una modifica
al punteggio non può cambiare e una al recupero sì.

\section{Le collezioni pubbliche: SciFact e NFCorpus}\label{sec:collezioni}

Due collezioni di BEIR~\cite{thakur2021beir}, scaricate con un controllo
dell'impronta degli archivi (tabella~\ref{tab:collezioni}). Servono a una cosa
sola: dire se il recupero lessicale di Koskidex ha un difetto, contro un
riferimento noto, BM25 di Lucene nella configurazione \emph{flat} di BEIR
(0,6789 e 0,3218 di nDCG@10~\cite{pyserini2cr}). Non rispondono alla domanda
della tesi, che vive sul corpus di dominio.

\begin{table}[ht]
  \centering
  \begin{tabular}{lrr}
    \toprule
    & SciFact & NFCorpus \\
    \midrule
    documenti & 5.183 & 3.633 \\
    query di test & 300 & 323 \\
    giudizi di test & 339 (1,1 per query) & 12.334 (38,2 per query) \\
    rilevanza & binaria & graduata (1 e 2) \\
    termini per query, mediana & 12 & 2 \\
    split di addestramento & train, 809 query & train 2.590, dev 324 \\
    \bottomrule
  \end{tabular}
  \caption{Le collezioni pubbliche.}
  \label{tab:collezioni}
\end{table}

\textbf{SciFact} è la prova di fumo: il riferimento è alto, e con un solo
documento pertinente per query una query si verifica a mano in un minuto.
\textbf{NFCorpus} è la prova vera: con 38 giudizi graduati per query il nDCG è
stabile, e la struttura somiglia più a un documentale, dove molti documenti sono
pertinenti in parte. Una trappola: il file
delle query contiene tutti gli split, e le query da valutare si ricavano dai
giudizi, altrimenti si misura su query che non ne hanno.

Sulle collezioni pubbliche Koskidex gira \emph{piatto}: in memoria, titolo e testo
in un campo solo, refusi spenti come nei riferimenti. La calibrazione, quando
c'è, usa gli split di addestramento e sviluppo, mai quello di test.

\section{Il corpus di dominio: gli albi pretori}\label{sec:corpus}

\paragraph{Perché gli albi.} Il corpus doveva venire dall'archivio di Documentale,
ma il progetto è stato abbandonato e non c'è nessun cliente a cui chiedere il
permesso di usarne i documenti. L'albo pretorio ha la forma di un documentale:
l'archivio di una sola organizzazione, con generi ricorrenti (determine, delibere,
ordinanze, avvisi, liquidazioni), un vocabolario suo e persone che cercano un atto
che sanno esistere. Ed è pubblico per obbligo di legge, quindi riscaricabile da
chiunque legga la tesi. Due fonti, verificate scaricandole il 23 settembre 2026:
\begin{itemize}
  \item \textbf{Comune di Crispiano} (TA), feed RSS, licenza CC BY 4.0: 563 atti
        da novembre 2025 a settembre 2026, 390 determine, 82 delibere, 58
        ordinanze. Il testo è nativo nei PDF, estraibile senza riconoscimento
        ottico: da due a sei pagine, fra 5 e 13 mila caratteri;
  \item \textbf{Regione Friuli Venezia Giulia}, API Socrata, licenza IODL 2.0:
        9.455 atti dal 2011, di 171 enti e 46 tipologie, con ente, tipologia,
        ufficio, oggetto e numero, ma senza testo integrale. I 171 enti si
        comportano come i clienti di un documentale condiviso.
\end{itemize}
Scartato AlboPOP, il progetto che converte gli albi in feed comuni: dei 215 feed
schedati ne rispondono 105, e scaricando a campione gli allegati ne arrivano 3 su
14. Gli altri rispondono \emph{200 OK} con una pagina d'errore: andava provato
scaricando, non contando le righe del catalogo.

\paragraph{Dentro Documentale.} Gli atti entrano in Documentale dal suo comando di
importazione, con il suo modello dei documenti, e in Elasticsearch dal suo comando
di indicizzazione, come in produzione. Il corpus per le valutazioni piatte esce
dall'esportazione di Documentale, non da una ricostruzione: un corpus costruito
scavalcando il sistema non direbbe niente sul sistema. L'indice ha 10.018 atti.

\paragraph{Dati personali.} Gli atti contengono nomi di persone: 79 oggetti del
Friuli Venezia Giulia hanno marcatori espliciti come \emph{nato a} o
\emph{codice fiscale}, e le pubblicazioni di
matrimonio di Crispiano riportano nomi, date e luoghi di nascita. Usarli per
misurare un motore è legittimo, ripubblicarli no: il corpus non entra nel
repository, che ne conserva solo gli script per riscaricarlo, e in questa tesi gli
esempi sono anonimi.

\section{Le query: known-item e bisogni aperti}\label{sec:query}

\paragraph{Known-item automatiche.} Chi sa il numero di un atto e il comune che
l'ha emesso lo ritrova? È la ricerca identificativa più comune in un documentale,
e si può fabbricare senza nessuno che annoti: la query nasce da un atto, e
quell'atto è l'unica risposta giusta. La forma è \texttt{<numero> <comune>}, per
esempio \texttt{13605 Fiume Veneto}. Degli atti con un numero che contiene una
cifra si tengono solo quelli la cui coppia (comune, numero) è unica: la
numerazione riparte ogni anno, e una query che descrive due atti non ha una sola
risposta. Su 9.806 candidati ne restano 9.425, e se ne estraggono 300 con un seme
fisso; altre 300, con un altro seme e nessun atto in comune, servono solo ad
addestrare e calibrare. Misurano una forma di query sola, e fabbricata: dicono se
il motore trova un atto di cui si sa il numero, non come lo cercherebbe una
persona.

\paragraph{Known-item umane.} Per le ricerche di un atto di cui si ricorda il
contenuto servono query scritte da persone. Il protocollo
(appendice~\ref{chapter:appendice-annotazione}) è fissato prima della raccolta: 160 atti
estratti con un seme, metà per fonte e dentro ogni fonte per genere, nessuno già
usato dalle known-item automatiche, divisi in quattro lotti da quaranta. A ogni
persona si mostra un atto come lo vedrebbe un impiegato, lo si nasconde, e le si
chiede cosa scriverebbe fra qualche settimana per ritrovarlo. L'atto si nasconde
prima che la persona scriva, altrimenti la query copia il titolo e la trova
qualunque motore. Scrivono persone che non conoscono i motori; l'autore di questa
tesi no, perché le sue query misurerebbero lui. Al momento di scrivere la raccolta
non è ancora fatta.

\paragraph{Bisogni aperti.} Per una query come \emph{manutenzione strade} non
esiste un atto giusto solo. Le 24 query del confronto, scritte a mano il 23
settembre 2026 in quattro famiglie (corte, medie, in lingua naturale, per numero
d'atto e con refusi), sono nate per vedere se i due motori trovano le stesse
cose, e sono servite a questo nel capitolo~\ref{chapter:sistemi}. Per misurarne la
pertinenza servono i giudizi della sezione seguente.

\section{Giudizi e accordo fra annotatori}\label{sec:giudizi}

\paragraph{Il pool.} Per le 24 query il pool è l'unione dei primi dieci risultati
di dodici configurazioni: le nove di Koskidex piatto misurate nei
capitoli~\ref{chapter:lessicale} e~\ref{chapter:ibrido}, dal baseline al solo
vettore, e tre passate dall'app, Elasticsearch di produzione e Koskidex innestato
con e senza vettori. Ne escono 905 righe da giudicare.

\paragraph{Il protocollo.} Prima di vedere un solo risultato, per ogni query si
scrive in una o due frasi cosa cerca chi la digita nel documentale di un comune,
e si committa: è l'unica difesa contro un bisogno che diventa «quello che i motori
hanno trovato». Poi si giudica ogni atto su tre gradi: 2 se risponde al bisogno
per intero, 1 se è pertinente ma parziale, 0 se non risponde, anche quando
contiene le parole della query. Le righe sono ordinate per identificativo, non per
posizione, e non dicono quale motore ha trovato l'atto; nel dubbio fra due gradi
vale il più basso, con una nota. Un giudizio non si cambia dopo aver visto le
metriche.

\paragraph{Il secondo annotatore.} Un solo annotatore misura anche il suo modo di
giudicare. Un secondo annotatore giudica, senza vedere i giudizi del primo,
query intere estratte con un seme dal foglio già annotato, almeno 150 righe. Si
riportano il kappa di Cohen sui tre gradi, il kappa pesato linearmente, perché
fra 1 e 2 si sbaglia meno che fra 0 e 2, e l'accordo semplice, comunque vengano.
Le metriche della tesi usano i giudizi del primo annotatore, e si dice di quanto
cambiano con quelli del secondo.

Anche i giudizi, come le known-item umane, sono da fare. Il
capitolo~\ref{chapter:confronto} li aspetta.

\paragraph{Le metriche.} Tre, calcolate dallo stesso codice per ogni motore:
\begin{itemize}
  \item \textbf{nDCG@10}, la metrica principale di BEIR, che pesa i gradi e la
        posizione nei primi dieci;
  \item \textbf{MRR@10}, il reciproco della posizione del primo documento
        pertinente, zero oltre il decimo: con un solo atto giusto, come nelle
        known-item, è la misura naturale;
  \item \textbf{Recall@100}, la quota dei pertinenti nei primi cento.
\end{itemize}
Accanto, sempre, le query a vuoto e il numero di candidati: una media non
distingue un motore che ordina male da uno che non trova, e i due chiedono
correzioni opposte.

\section{Gli strumenti}\label{sec:strumenti}

\paragraph{Koskidex piatto.} \texttt{scripts/evaluate} carica una collezione in
formato BEIR in memoria, la indicizza con le impostazioni date da riga di comando,
esegue le query e scrive metriche, candidati e tempi per query. Con i vettori
calcola gli embedding una volta e li tiene in una cache.

\paragraph{Dall'app.} \texttt{app:eval-run-queries}, un comando aggiunto a
Documentale, esegue un file di query con la ricerca di produzione, sul motore
scelto dalla configurazione, e scrive per ogni query gli identificativi trovati e
il tempo della chiamata. Lo stesso \texttt{scripts/evaluate}, con un'opzione, ne
calcola le metriche: i motori passati dall'app e Koskidex piatto si contano con lo
stesso codice.

\paragraph{Confronto e pool.} \texttt{scripts/compare} esegue le stesse query su
più configurazioni di Koskidex in memoria; \texttt{scripts/pool} costruisce il
foglio dei giudizi dalle configurazioni di Koskidex e dai rapporti dell'app.

\paragraph{Raccolta e giudizi.} Le query umane e i giudizi non passano da un
foglio di calcolo. Due pagine HTML, da aprire nel browser senza rete, fanno da
sole i passi dei protocolli: mostrano un atto alla volta, lo nascondono quando
serve, applicano l'ordine e nascondono il motore, registrano i tempi e le
sessioni, e alla fine scaricano un file da rimandare. La pagina dei giudizi si
rifiuta di esistere finché i bisogni non sono tutti scritti e committati. Due
script importano i file ricevuti, controllano che corrispondano alla pagina da cui
vengono, e scrivono le query e i giudizi in formato BEIR. Le pagine contengono il
testo degli atti e non entrano nel repository.

\paragraph{La macchina.} Tutti i tempi sono misurati su un Apple M2 con 8 core e
16 GB. Koskidex gira nativo su tutti i core; Elasticsearch 9.1.0 e MySQL in una
macchina virtuale Docker con 4 CPU, Elasticsearch con 1 GB di heap, e il suo
tempo include la chiamata HTTP da PHP. I tempi dei due motori non si confrontano
alla pari, e dove stanno fianco a fianco lo si scrive.
